\scanpage{22}
\mspage{16}

\textcircled{4} Les hypothèses préliminaires sur $(M,W)$, y
\inserted{\unsure{inclus|incluse} l'accessibilité,}
ainsi que les conditions équivalentes du \textcircled{2},
sont stables par passage de $(M,W)$ à
$(M(I),W(I))$. Mais qu'en est-il de la
condition du \textcircled{3} ? Est-elle conséquence
des conditions du \textcircled{2} (i.e. leur est-elle
équivalente ?).

Je vais esquisser la démonstration du
théorème\inserted{-scholie}, pour cerner le travail technique
qui reste à faire.

$1^\circ$) est de la vérification de routine.

$2^\circ$) Voici les implications $\pm$ tautologiques
entre les conditions envisagées :
\[
 \text{les}\qquad a''\Rightarrow a'\overset{?}{\Rightarrow}a\Rightarrow b\deleted{\unsure{\Rightarrow b'}}\Leftarrow b',
\]
mais je vois que $a'\Rightarrow a$ n'est \deleted{pas} nullement
``tautologique'', et peut-être même faux
après tout !\textsuperscript{$*$} Il faudrait supposer que dans
% slanted
\margin{\unsure{C'est pourquoi} j'ai modifié \unsure{$a'$)}, \unsure{en imposant|en renforçant} à a$'$) \deleted{\unsure{pour}} \unsure{l'égalité des|j'y ai mis} \unsure{?} après le trait !}
le quadruplet $(C,F,TC,TF)$, $F$ et $TF$ soient clos
à droite, \inserted{i.e. $\unsure{F=}\deleted{\unsure{TF}}$ et \uline{plus précisément} que $F=(TC)_*$, $TF=\unsure{,}C_*$,} ce qui sera le cas p.ex. si
on a une structure de $Q$-catégorie de
modèles ``close'' (``closed model category''), au
sens de Quillen. Les autres \inserted{trois} implications dans
\unsure{(*)} sont bel et bien tautologiques.

$a\Rightarrow a'$ On peut dans a) supposer $\deleted{\unsure{TC_0}}$
\[
 TC_0\subset C_0
\]
(quitte à remplacer $C_0$ par $C_0\cup TC_0$), puis, grâce à
% archive sheet number: 21
